\documentclass{beamer}
\mode<presentation> {%
	\usetheme[]{Boadilla}
	\usecolortheme{orchid}
	\setbeamercovered{transparent}
	\setbeamersize{text margin left=0.7cm}
  \setbeamersize{text margin right=0.7cm}
  \setbeamertemplate{enumerate item}[square]
  \setbeamertemplate{enumerate subitem}[circle]
  \setbeamertemplate{enumerate subsubitem}[>]
  \setbeamertemplate{itemize item}[square]
  \setbeamertemplate{itemize subitem}[circle]
  \setbeamertemplate{itemize subsubitem}[>]
}

\usepackage[utf8]{inputenc}
\usepackage{graphicx}


\title[Procesamiento de Señales]{Caracterización en frecuencia de señales}
\author[J. Chacón]{%[Author, Another] % (optional, use only with lots of authors)
\textit{Ingeniería Electrónica de Comunicaciones }\\
\vspace{0.6cm}
\small Jesús Chacón Sombría}

\institute[DACYA] % (optional, but mostly needed)
{
  %\inst{1}%

\footnotesize  Departamento de Arquitectura de Computadores y Automática.\\
  Universidad Complutense de Madrid
  }

\date[Transformada Fourier Discreta] % (optional)
{\footnotesize Curso 2024-2025}%{\today / Master}
\subject{Talks}

\pgfdeclareimage[height=0.75cm]{logo}{ucm.png}

\usepackage[absolute,overlay]{textpos}
\setlength{\TPHorizModule}{1mm}
\setlength{\TPVertModule}{1mm}
\newcommand{\MyLogo}{%
\begin{textblock}{14}(116.0,0.6)
  \pgfuseimage{logo}
\end{textblock}
}

\makeatletter
\numberwithin{section}{part}
\AtBeginPart{
\beamer@tocsectionnumber=0\relax
\frame{\partpage}
}
\makeatother

\begin{document}

\frame{\titlepage}

\section{Introducción}
\begin{frame}{Introducción}
    \begin{itemize}
        \item Las señales en distintos dominios tienen contenido espectral característico.
        \item Comprender la distribución de frecuencias es clave para el procesamiento y análisis de señales.
    \end{itemize}
\end{frame}

\section{Comunicaciones}
\begin{frame}{Comunicaciones}
    \begin{itemize}
        \item \textbf{Señal:}
        \begin{itemize}
            \item \textbf{Señales de voz:} 300 Hz – 3.4 kHz (telefonía), hasta 20 kHz en alta fidelidad.
            \item \textbf{Radio:}
            \begin{itemize}
                \item AM: 300 kHz – 3 MHz
                \item FM: 88 – 108 MHz
                \item Wi-Fi: 2.4 GHz y 5 GHz
            \end{itemize}
            \item \textbf{Fibra óptica:} decenas a centenas de THz.
        \end{itemize}
        \item \textbf{Fuentes de ruido:}
        \begin{itemize}
            \item interferencias electromagnéticas (EMI) ($>$ kHz – GHz),
            \item ruido térmico (afecta toda la banda),
            \item desvanecimiento de señales (Hz – MHz).
        \end{itemize}
    \end{itemize}
\end{frame}

\section{Audio y Música}
\begin{frame}{Audio y Música}
    \begin{itemize}
        \item \textbf{Señal:}
        \begin{itemize}
            \item Espectro audible: 20 Hz – 20 kHz.
            \item Bandas:
            \begin{itemize}
                \item Subgraves: 20 – 60 Hz
                \item Graves: 60 – 250 Hz
                \item Medios: 250 Hz – 2 kHz
                \item Agudos: 2 kHz – 20 kHz
            \end{itemize}
            \item Frecuencias de muestreo típicas:\\8 kHz, 44.1 kHz, 48 kHz, 192 kHz.
        \end{itemize}
        \item \textbf{Fuentes de ruido:}
        \begin{itemize}
            \item ruido de cuantización ($>$10 kHz),
            \item interferencias eléctricas (50/60 Hz y armónicos),
            \item ruido ambiental (banda ancha 20 Hz – 20 kHz).
        \end{itemize}
    \end{itemize}
\end{frame}

\section{Señales Biomédicas}
\begin{frame}{Señales Biomédicas}
    \begin{itemize}
        \item \textbf{Señal}
        \begin{itemize}\scriptsize
            \item \textbf{ECG (Electrocardiograma):} 0.05 – 100 Hz.\\Utilizado para analizar la actividad eléctrica del corazón.
            \item \textbf{EEG (Electroencefalograma):} 0.1 – 100 Hz.\\Clasificación de ondas cerebrales:
            \begin{itemize}\scriptsize
                \item Delta: 0.1 – 4 Hz (sueño profundo)
                \item Theta: 4 – 8 Hz (relajación ligera)
                \item Alfa: 8 – 14 Hz (estado de calma)
                \item Beta: 14 – 30 Hz (actividad mental)
                \item Gamma: $>$30 Hz (procesos cognitivos altos)
            \end{itemize}
            \item \textbf{EMG (Electromiografía):} 10 Hz – 500 Hz.\\Registra la actividad eléctrica de los músculos.
            \item \textbf{PCG (Fonocardiografía):} 10 Hz – 1 kHz.\\Se usa para el análisis de sonidos cardíacos.
            \item \textbf{Ultrasonido médico:} 1 – 15 MHz.\\Utilizado en imágenes de diagnóstico como ecografías.
        \end{itemize}
        \item \textbf{Fuentes de ruido:}
        \begin{itemize}\scriptsize
            \item interferencias eléctricas (50/60 Hz),
            \item artefactos de movimiento ($<$10 Hz),
            \item ruido fisiológico (baja frecuencia $<$1 Hz – 100 Hz).
        \end{itemize}
    \end{itemize}
\end{frame}

\section{Sistemas de Control y Robótica}
\begin{frame}{Sistemas de Control y Robótica}
    \begin{itemize}
        \item \textbf{Señal:}
        \begin{itemize}
            \item Sensores industriales: $<$ 1 kHz.
            \item Visión artificial: 30 – 60 FPS (Nyquist de 15 – 30 Hz).
            \item Robots móviles y manipuladores:
            \begin{itemize}
                \item Sensores de posición y velocidad: $<$1 kHz.
                \item Sensores de fuerza y torque: $<$500 Hz.
                \item Cámaras y LiDAR: 10 Hz – 1 kHz.
            \end{itemize}
            \item Control de motores eléctricos:
            \begin{itemize}
                \item Frecuencias de conmutación PWM: 1 kHz – 50 kHz.
                \item Frecuencias de vibración y resonancia: 10 Hz – 10 kHz.
            \end{itemize}
            \item Redes de comunicación en robótica:
            \begin{itemize}
                \item CAN bus: 125 kHz – 1 MHz.
                \item EtherCAT y protocolos industriales: MHz.
            \end{itemize}
        \end{itemize}
        \item \textbf{Fuentes de ruido:}
        \begin{itemize}
            \item variaciones en sensores (baja frecuencia $<$1 kHz),
            \item errores de cuantización ($>$Nyquist),
            \item interferencias electromagnéticas ($>$kHz),
            \item vibraciones mecánicas (10 Hz – 10 kHz),
            \item ruido térmico en sensores ($<$100 Hz),
        \end{itemize}
    \end{itemize}
\end{frame}


\begin{frame}{Sistemas de Control y Robótica}
    \begin{itemize}
        \item \textbf{Señal:}
        \begin{itemize}
            \item Sensores de localización y navegación:
            \begin{itemize}
                \item \textbf{GPS:} 1.1 – 1.6 GHz.
                \item \textbf{IMU (Inertial Measurement Unit):} 1 Hz – 1 kHz\\frecuencias típicas en acelerómetros y giróscopos.
                \item \textbf{Sonar:} 10 kHz – 1 MHz dependiendo de la aplicación\\ (baja frecuencia para largo alcance, alta frecuencia para mayor resolución).
                \item \textbf{LiDAR:} pulsos láser en 10 Hz – 1 kHz.
                \item \textbf{UWB (Ultra-Wideband):} 3.1 – 10.6 GHz\\utilizado para localización de alta precisión.
            \end{itemize}
        \end{itemize}
        \item \textbf{Fuentes de ruido:}
        \begin{itemize}
            \item variaciones en sensores (baja frecuencia $<$1 kHz),
            \item interferencias electromagnéticas ($>$kHz),
            \item vibraciones mecánicas (10 Hz – 10 kHz),
            \item ruido térmico en sensores (<100 Hz),
            \item interferencias multipath en GPS (MHz – GHz).
        \end{itemize}
    \end{itemize}
\end{frame}


\section{Geofísica y Sismología}
\begin{frame}{Geofísica y Sismología}
    \begin{itemize}
        \item \textbf{Señales:}
        \begin{itemize}\scriptsize
            \item \textbf{Terremotos:} $<$ 1 Hz – 10 Hz.\\Contienen componentes de baja frecuencia que pueden propagarse a largas distancias.
            \item \textbf{Microtremores:} 0.1 – 1 Hz.\\Se generan por actividad humana y fenómenos naturales menores.
            \item \textbf{Vibraciones en estructuras:}\\0.1 Hz – 100 Hz. Importantes para la detección de resonancias y evaluación de estabilidad.
            \item \textbf{Ondas sísmicas:}
            \begin{itemize}\scriptsize
                \item Ondas P (primarias): Frecuencia entre 0.1 – 10 Hz.\\Son las más rápidas y viajan a través de sólidos y líquidos.
                \item Ondas S (secundarias): 0.1 – 3 Hz.\\Se propagan solo en sólidos y son más destructivas que las ondas P.
                \item Ondas superficiales (Rayleigh y Love): 0.01 – 1 Hz.\\Son las responsables del daño en estructuras.
            \end{itemize}
        \end{itemize}
        \item \textbf{Fuentes de ruido:} 
        \begin{itemize}\scriptsize
            \item actividad humana (tráfico, maquinaria) ($<$10 Hz),
            \item ruido atmosférico ($<$1 Hz),
            \item interferencias ambientales (baja frecuencia $<$10 Hz).
        \end{itemize}
    \end{itemize}
\end{frame}

\end{document}
